\documentclass[11pt]{article}
\usepackage{ochamath}
\subjectcolor{2F5DA8}
\setsheet{線形代数・電気系院試補充}{LU・QR・擬似逆行列　演習}{https://university-math-crj.pages.dev/linear-algebra/rank-factorizations/}
\namelinetrue
\begin{document}
\makesheethead{問題}
自作問題。本文の例題とは数値や設定が異なります。条件・途中式・検算を示してください。
\problem[LUと連立式]
$A=\begin{pmatrix}4&1\\2&3\end{pmatrix},b=(6,8)$ をLU分解で解け。
\workspace{45mm}
\problem[QR最小二乗]
$A=\begin{pmatrix}1&0\\1&1\\0&1\end{pmatrix},b=(2,1,0)$ をQRで最小二乗近似せよ。
\workspace{45mm}
\newpage
\problem[階数1の擬似逆]
$A=\begin{pmatrix}2&2\\1&1\end{pmatrix},b=(0,1)$。擬似逆と最小ノルム最小二乗解を求めよ。
\workspace{45mm}
\problem[条件数]
$A=\operatorname{diag}(2,0.02)$ の $\kappa_2(A),\kappa_2(A^{\mathsf T}A)$ を求めよ。
\workspace{45mm}
\newpage
\problem[Moore--Penrose条件]
$A=\operatorname{diag}(3,0)$ の $A^+$ を求め、4つの条件を確認せよ。
\workspace{45mm}
\problem[列ベクトル化]
$A=\begin{pmatrix}1&2\\0&3\end{pmatrix},B=\operatorname{diag}(2,4),X=\begin{pmatrix}1&3\\2&4\end{pmatrix}$ で vec の公式を確かめよ。
\workspace{45mm}
\end{document}
